\section{Chương~\ref{sec:code}. Mã Morse}

Các ước tính giới hạn của chương là lý thuyết thông tin và phép tính; những gì đã đo là nhiễu phát sinh ở đâu trong kênh, não làm việc nhanh đến đâu và một xung tốn bao nhiêu, còn câu hỏi vỏ não thực sự dùng mã nào vẫn để ngỏ.

{\footnotesize
\setlength{\tabcolsep}{3pt}\renewcommand{\arraystretch}{1.15}
\begin{longtable}{>{\raggedright}p{0.5cm}>{\raggedright}p{4.0cm}>{\raggedright}p{5.6cm}>{\raggedright}p{2.3cm}>{\raggedright\arraybackslash}p{1.7cm}}
\toprule
TT & Khẳng định & Thí nghiệm và điều đã đo & Nguồn & Trạng thái \\
\midrule\endfirsthead
\toprule
TT & Khẳng định & Thí nghiệm và điều đã đo & Nguồn & Trạng thái \\
\midrule\endhead
1 & Tần số của nơ-ron \nt{GLUT} vỏ não từ một phần nhỏ đến hàng chục xung mỗi giây, và phần lớn thời gian nơ-ron làm việc gần giới hạn dưới & Ghi bằng pipet áp sát tế bào (việc chọn nơ-ron không phụ thuộc vào hoạt động của nó) ở vỏ thính giác chuột cống đang thức: tại mỗi thời điểm, dưới $5\,\%$ nơ-ron có tần số trên $20$ xung mỗi giây; một phần nơ-ron có tần số nền dưới $0.01$ mỗi giây; phân bố tần số là log-chuẩn & \cite{Hromadka2008} & đã đo \\
2 & Dung lượng của kênh xung $R_{\max}\approx r\log_2\frac{e}{r\Delta t}$~\eqref{eq:spikeentropy}, gần như toàn bộ nằm ở thời điểm xung (mệnh đề~\ref{prop:capacity}) & Entropy của chuỗi nhị phân gồm các ô dài $\Delta t$. Các con số của chương: $8$ bit trên một xung khi $r=10$ mỗi giây và $\Delta t=1\ms$, khoảng $5$ bit khi $\Delta t=10\ms$, dưới $2$ bit mỗi giây với bộ đếm xung & \cite{MacKay1952} & lý thuyết \\
3 & Nơ-ron cảm giác truyền từ một đến vài bit trên một xung, trong trường hợp tốt nhất tới một nửa giới hạn & Khôi phục kích thích từ xung của các nơ-ron cảm giác mà kích thích đã biết; tổng hợp các phép đo trong sách & \cite{Rieke1997} & đã đo \\
4 & Bộ tạo xung chính xác; thời điểm chính xác do những biến đổi nhanh của đầu vào quyết định & Nơ-ron vỏ não chuột cống trong lát cắt: với dòng lặp lại có dao động nhanh, xung lặp lại với độ chính xác dưới $1\ms$; với dòng không đổi, thời điểm trôi dạt & \cite{Mainen1995} & đã đo \\
5 & Khớp thần kinh không đáng tin cậy: hơn một nửa số xung không tới được nơi nhận do tính ngẫu nhiên của việc nhả & Kích thích tối thiểu các đầu vào của nơ-ron tháp hồi hải mã (lát cắt): hơn một nửa số xung không gây đáp ứng. Đã loại trừ dao động ngưỡng kích thích, lỗi dẫn truyền ở chỗ phân nhánh sợi trục và nhiệt độ thí nghiệm thấp; còn lại là sự nhả mang tính xác suất & \cite{AllenStevens1994} & đã đo \\
6 & Dòng xung trong vỏ não sống trông ngẫu nhiên: các khoảng phân tán như ở dòng Poisson, phương sai số xung gần bằng trung bình; một phần ``nhiễu'' là thông điệp về cử động của con vật & Ghi ở các vùng thị giác V1 và MT của khỉ đang thức: hệ số biến thiên của các khoảng xấp xỉ $1$, tỷ số phương sai số xung trên trung bình như ở một quá trình ngẫu nhiên; mô hình cộng các đầu vào nhỏ độc lập cho độ phân tán nhỏ hơn nhiều. Ở vỏ thị giác chuột nhắt, một phần đáng kể độ phân tán dự đoán được từ video ghi cử động của chính con vật & \cite{Softky1993,Stringer2019} & đã đo \\
7 & Mã tần số chậm: sai số $1/\sqrt{rT}$~\eqref{eq:rateerror}, còn não nhận ra hình ảnh trong $\sim150\ms$ & Công thức là thống kê Poisson, suy luận của chương. Thực nghiệm: người tham gia quyết định có con vật trong một bức ảnh lạ được chiếu $20\ms$ hay không; điện thế não với ``có'' và ``không'' tách nhau sau khoảng $150\ms$. Việc chia thời gian này thành khoảng một chục tầng, mỗi tầng $10$--$20\ms$, là ước tính của chương, không phải phép đo & \cite{Thorpe1996} & đã đo \\
8 & Lấy trung bình trên nhóm bị tương quan giới hạn: sai số giảm không quá $1/\sqrt c$ lần~\eqref{eq:correlated} (mệnh đề~\ref{prop:pooling}) & Các cặp nơ-ron vùng MT của khỉ được ghi đồng thời: hệ số tương quan số xung trung bình $0.12$, và phân tích lý thuyết cho thấy ngay cả tương quan như vậy cũng giới hạn lợi ích của nhóm. Lý thuyết: chỉ phần tương quan giống tín hiệu mới đặt ra giới hạn. Mô hình của quan điểm đối lập: tần số của nhóm $50$--$100$ nơ-ron được đọc trong một khoảng giữa hai xung, $10$--$50\ms$, và thời điểm của từng xung không được vỏ não sử dụng & \cite{Zohary1994,MorenoBote2014,ShadlenNewsome1998} & đã đo \\
9 & Con số có thể được ghi vào thời điểm xung đầu tiên~\eqref{eq:latency}, vào thứ tự xung của một nhóm hoặc vào pha của nhịp cục bộ & Võng mạc kỳ giông: cấu trúc không gian của một hình ảnh được chiếu ngắn được ghi trong độ trễ tương đối của các xung đầu tiên của nơ-ron đầu ra, gần như không phụ thuộc độ tương phản. Tế bào vị trí ở hồi hải mã chuột cống: khi đi qua vị trí ``của mình'', xung dịch ngày càng sớm hơn trong chu kỳ nhịp theta ($7$--$12$\,Hz), độ dịch từ $100^\circ$ đến $355^\circ$. Vỏ não dùng thời điểm xung đến mức nào là câu hỏi mở (dòng~8) & \cite{Gollisch2008,OKeefe1993} & đã đo \\
10 & Mã nào được đọc do hằng số thời gian của nơi nhận quyết định~\eqref{eq:reader} (mệnh đề~\ref{prop:reader}); ở vỏ não đang hoạt động, nơ-ron gần với bộ phát hiện trùng khớp hơn & Công thức~\eqref{eq:reader} là suy luận của chương. Tổng quan các bản ghi in vivo: ở vỏ não đang hoạt động, độ dẫn của màng cao hơn vài lần so với lúc nghỉ, và hằng số thời gian ngắn hơn tương ứng. Chưa có bằng chứng trực tiếp rằng vỏ não chuyển mã đúng theo cách này & \cite{Destexhe2003} & giả thuyết \\
11 & Khớp thần kinh cạn kiệt truyền sự thay đổi tần số, về thực chất là $\Delta r/r$~\eqref{eq:depression} (hình~\ref{fig:depression}) & Ghi cặp nơ-ron tháp vỏ não: tốc độ cạn kiệt do xác suất nhả quy định; khi cạn kiệt chậm, đáp ứng phản ánh tần số, khi nhanh --- phản ánh sự trùng khớp. Mô hình dựa trên sự cạn kiệt đo được ở vỏ não chuột cống: những thay đổi tương đối như nhau của tần số ở đầu vào nhanh và chậm cho đáp ứng như nhau, tức là khớp thần kinh truyền $\Delta r/r$. Các con số $1.7\to2.2$, $6.7$ và $90\ms$ là tính toán của chương & \cite{Tsodyks1997,Abbott1997depression} & đã đo \\
12 & Chùm xung là ``gạch'': xác suất nó bị mất $(1-p)^k$~\eqref{eq:burst} là nhỏ, tạo thuận còn làm giảm thêm & Tổng quan: ở các khớp thần kinh không tin cậy, xung đơn thường bị mất, còn chùm xung được truyền đáng tin cậy nhờ tạo thuận; chùm xung gây ra những thay đổi lâu dài của khớp thần kinh. Việc não đọc chùm xung như một ký hiệu riêng là giả thuyết & \cite{Lisman1997,AllenStevens1994} & giả thuyết \\
13 & Nơ-ron \nt{GABA} nhanh tạo nhịp và ``dấu câu''; một lớp khác ức chế các nơ-ron ức chế và mở cổng (mệnh đề~\ref{prop:punctuation}) & Tổng quan tính chất các lớp nơ-ron \nt{GABA} vỏ não. Quang di truyền: kích thích nhịp nhàng các nơ-ron \nt{GABA} nhanh gây ra nhịp gamma, và đáp ứng với kích thích phụ thuộc pha của chu kỳ. Ở vỏ thị giác chuột nhắt, nơ-ron PV ức chế lẫn nhau, nơ-ron SST ức chế mọi lớp khác, nơ-ron VIP chủ yếu ức chế SST. Kích thích nơ-ron VIP ở chuột nhắt đang thức gỡ bỏ ức chế khỏi nơ-ron \nt{GLUT}; chúng được bật bởi tín hiệu củng cố. Sự phân vai như một quy tắc chung là giả thuyết & \cite{Tremblay2016,Cardin2009,Pfeffer2013,Pi2013} & giả thuyết \\
14 & Năng lượng giới hạn tần số trung bình ở mức cỡ một xung mỗi giây~\eqref{eq:energy} & Ngân sách năng lượng của chất xám loài gặm nhấm theo dữ liệu giải phẫu và sinh lý: xung và dòng qua khớp thần kinh chiếm $47\,\%$ và $34\,\%$ năng lượng truyền tín hiệu; cùng lúc chỉ có thể có không quá $\sim15\,\%$ nơ-ron hoạt động. Với vỏ não người: không quá $\sim1\,\%$ nơ-ron có thể hoạt động mạnh cùng lúc. Con số $\sim10^9$ ATP mỗi xung và công suất $\sim10$\,W cho truyền tín hiệu là ước tính của chương dựa trên các tính toán này & \cite{Attwell2001,Lennie2003} & lý thuyết \\
15 & Mã thưa và được tinh chỉnh cho số bit trên mỗi joule lớn nhất; vỏ não đổi độ chính xác lấy năng lượng (mệnh đề~\ref{prop:economy}) & Giả thuyết mã tiết kiệm. Cơ sở: ở chuột nhắt bị hạn chế thức ăn, độ dẫn của thụ thể AMPA và mức tiêu thụ ATP của khớp thần kinh ở vỏ thị giác giảm ($-29\,\%$), tần số xung được giữ nguyên, còn độ chọn lọc hướng nới rộng $32\,\%$ & \cite{Barlow1961,Padamsey2022} & giả thuyết \\
\bottomrule
\end{longtable}}
